\documentclass[12pt,a4paper]{article}
\usepackage[utf8]{inputenc}
\usepackage[T2A]{fontenc}
\usepackage[russian]{babel}
\usepackage{amsmath}
\usepackage{amsfonts}
\usepackage{amssymb}
\usepackage{geometry}
\usepackage{tensor}
\geometry{left=2cm,right=2cm,top=2cm,bottom=2cm}

\title{\textbf{В рамочку обведено то, что скорее всего нужно написать в тестике}}
\author{}
\date{}

\begin{document}

\maketitle

\section*{Задача 1. Потенциал диполя и квадруполя}

\textbf{Условие:} Записать выражения для потенциала диполя и квадруполя.

\textbf{Решение:}

Разложение потенциала в мультипольный ряд:

\begin{equation}
\varphi(\vec{r}) = \int d^3r' \frac{\rho(\vec{r}')}{|\vec{r} - \vec{r}'|}
\end{equation}

При $r \gg r'$ разлагаем:

\begin{equation}
\frac{1}{|\vec{r} - \vec{r}'|} = \frac{1}{r} + \frac{\vec{r} \cdot \vec{r}'}{r^3} + \frac{3(\vec{r} \cdot \vec{r}')^2 - r^2 r'^2}{2r^5} + \ldots
\end{equation}

\textbf{Дипольный член:}
\begin{equation}
\boxed{\varphi^{(1)}(\vec{r}) = \frac{\vec{d} \cdot \vec{r}}{r^3}}
\end{equation}
где дипольный момент:
\begin{equation}
\boxed{\vec{d} = \int d^3\vec{r}' \, \vec{r}' \, \rho(\vec{r}')}
\end{equation}

\textbf{Квадрупольный член:}
\begin{equation}
\boxed{\varphi^{(2)}(\vec{r}) = \frac{D_{ij} r_i r_j}{2r^5}}
\end{equation}
где тензор квадрупольного момента:
\begin{equation}
\boxed{D_{ij} = \int d^3\vec{r}' \, \rho(\vec{r}') \left(3r'_i r'_j - \delta_{ij} r'^2\right)}
\end{equation}

\section*{Задача 2. Инвариантность дипольного момента}

\textbf{Условие:} Показать, что для электронейтральной системы зарядов дипольный момент не зависит от выбора начала координат.

\textbf{Решение:}

Пусть система зарядов электронейтральна: $\sum_\alpha e_\alpha = 0$.

Дипольный момент относительно начала координат $O$:
\begin{equation}
\boxed{\vec{d} = \sum_\alpha e_\alpha \vec{r}_\alpha}
\end{equation}

При смещении начала координат на $\vec{r}_0$:
\begin{equation}
\boxed{\vec{d}' = \sum_\alpha e_\alpha (\vec{r}_\alpha - \vec{r}_0) = \sum_\alpha e_\alpha \vec{r}_\alpha - \vec{r}_0 \sum_\alpha e_\alpha}
\end{equation}

Так как $\sum_\alpha e_\alpha = 0$:
\begin{equation}
\boxed{\vec{d}' = \vec{d}}
\end{equation}

Дипольный момент инвариантен относительно выбора начала координат для электронейтральной системы.

\section*{Задача 3. Потенциал от плотности заряда}

\textbf{Условие:} Записать выражение для потенциала, создаваемого распределением заряда с плотностью $\rho(\vec{r})$.

\textbf{Решение:}

Решение уравнения Пуассона $\Delta \varphi = -4\pi \rho$:

\begin{equation}
\boxed{\varphi(\vec{r}) = \int d^3\vec{r}' \frac{\rho(\vec{r}')}{|\vec{r} - \vec{r}'|}}
\end{equation}

Это частное решение неоднородного уравнения. Общее решение включает также решение однородного уравнения $\Delta \varphi_0 = 0$.

\section*{Задача 4. Волновые уравнения в калибровке Лоренца}

\textbf{Условие калибровки Лоренца:}
\[
\frac{1}{c}\frac{\partial \varphi}{\partial t} + \operatorname{div}\vec{A} = 0 \quad \left(\text{или } \boxed{\partial_\mu A^\mu = 0}\right)
\]

\textbf{Волновые уравнения для потенциалов:}
\[
\Box \varphi = 4\pi\rho \quad \Longrightarrow \quad \boxed{\frac{1}{c^2}\frac{\partial^2 \varphi}{\partial t^2} - \Delta \varphi = 4\pi\rho}
\]
\[
\Box \vec{A} = \frac{4\pi}{c}\vec{j} \quad \Longrightarrow \quad \boxed{\frac{1}{c^2}\frac{\partial^2 \vec{A}}{\partial t^2} - \Delta \vec{A} = \frac{4\pi}{c}\vec{j}}
\]
где $\Box = \dfrac{1}{c^2}\dfrac{\partial^2}{\partial t^2} - \Delta$ --- оператор Даламбера.

\textbf{В ковариантной форме (единым уравнением для 4-потенциала $A^\mu$):}
\[
\boxed{\Box A^\mu = \frac{4\pi}{c} j^\mu}
\]
\section*{Задача 5. II пара уравнений Максвелла и проверка согласованности}

\textbf{Условие:} Записать II пару уравнений Максвелла и проверить их согласованность с уравнением непрерывности.

\textbf{Решение:}

\textbf{II пара уравнений Максвелла} в ковариантной форме:
\begin{equation}
\boxed{\frac{\partial F^{\mu\nu}}{\partial x^\nu} = -\frac{4\pi}{c} j^\mu}
\end{equation}

В 3-мерной форме:
\begin{equation}
\boxed{\text{div} \vec{E} = 4\pi \rho}
\end{equation}
\begin{equation}
\boxed{\text{rot} \vec{H} = \frac{1}{c}\frac{\partial \vec{E}}{\partial t} + \frac{4\pi}{c} \vec{j}}
\end{equation}

\textbf{Проверка согласованности:}

Берём дивергенцию от второго уравнения:
\begin{equation}
\nabla \cdot (\nabla \times \vec{H}) = \frac{1}{c}\frac{\partial}{\partial t}(\nabla \cdot \vec{E}) + \frac{4\pi}{c} \nabla \cdot \vec{j}
\end{equation}

Так как $\nabla \cdot (\nabla \times \vec{H}) = 0$ и $\nabla \cdot \vec{E} = 4\pi \rho$:
\begin{equation}
0 = \frac{4\pi}{c}\frac{\partial \rho}{\partial t} + \frac{4\pi}{c} \nabla \cdot \vec{j}
\end{equation}

Получаем \textbf{уравнение непрерывности}:
\begin{equation}
\boxed{\frac{\partial \rho}{\partial t} + \nabla \cdot \vec{j} = 0}
\end{equation}

В ковариантной форме: $\partial_\mu j^\mu = 0$.

Согласованность доказана: II пара уравнений Максвелла автоматически удовлетворяет закону сохранения заряда.

\end{document}